\chapter{Scope and current state}

This project formalizes the characterization of fibers of the principal minor map.
For a finite index type $I$ and a field $F$, write $A_S=\det A[S,S]$ and
\[ f_A=\det(\operatorname{diag}(x_i)_{i\in I}+A),\qquad
   G_A=\operatorname{adj}(\operatorname{diag}(x_i)+A). \]
The question is when the principal minors of $A$ determine it up to invertible diagonal
conjugation and transpose, that is, up to diagonal equivalence.

There are two targets.
\begin{enumerate}
\item Al Ahmadieh, \emph{The fiber of the principal minor map} (arXiv:2309.00806),
Theorem 4.1.
\item Chatterjee--Ghosh--Gurjar--Raj, \emph{Characterizing and Testing Principal Minor
Equivalence of Matrices} (STOC 2025, arXiv:2410.01961), Theorem 1.1
(Theorem~\ref{thm:cggr11}): irreducible matrices have the same principal minors exactly
when they are related by a sequence of cut-transposes followed by a diagonal equivalence.
The alternate Theorem 4.1 is a corollary. The algorithmic results of that paper are out of
scope.
\end{enumerate}

\section{Libraries}

\texttt{DetPoly} is the library. It contains the definitions, diagonal similarity and
equivalence, irreducible matrices, the Desnanot--Jacobi identity, Jacobi's formula for a
derivation, and Theorem 2.3, all proved (Chapter~\ref{chap:detpoly}). Two statements toward
Lemma 4.3 in the library's own terms are unfinished. It also contains two self-contained
developments on the index type $\{0,\dots,n-1\}$: a complete proof of Lemma 4.3 given
Lemma 4.2, and of the shuffle sign of Lemma 3.2.

\texttt{LEAN\_Warmup} is exploratory. It contains the draft statements for the paper,
the outline of CGGR with the alternate theorem, generalized Laplace expansion, tools for
multivariate polynomials, a computable polynomial model, and problem-set proofs.
Warmup modules may import the library; the library does not import the warmups.
\texttt{LEAN\_Warmup.JoeWarmup} exports no declarations, and the diagonal expansion of
Theorem~\ref{thm:diagonal-expansion} currently exists in both libraries under the same
names.

\section{Reading the dependency graph}

Every explicitly named public definition and result in both roots is linked from some
node, including helper APIs and examples; grouped nodes link all declarations implementing
one piece of mathematics. A Lean link records that a statement or definition exists.
A proof marked complete means that the Lean proof of that node contains no \texttt{sorry}
of its own; it may still depend on unfinished nodes, which the graph shows.
A node without a proof carries a status line: \emph{to do} or \emph{unfinished} for our own
proofs, \emph{external} for results cited from other papers. Dependency arrows describe
mathematical prerequisites or the intended proof route for unfinished nodes; they are not
a literal Lean import graph.

\section{Dependencies}

The Lean proof of Theorem~\ref{thm:main-alt} is complete given:
\begin{itemize}
\item cut-transpose preserves principal minors (Lemma~\ref{lem:ct-pme}), the first real
proof obligation;
\item reducible matrices have nontrivial fibers (Lemma~\ref{lem:reducible-fiber});
\item irreducibility of $A$ is irreducibility of $f_A$ (Lemma~\ref{lem:irred-poly}), which
uses the external Lemma~\ref{ext:nonzero};
\item equality of fibers is equality of determinantal polynomials
(Lemma~\ref{lem:fiber-poly});
\item compatibility of cut-transposes and cuts with diagonal equivalence
(Lemmas~\ref{lem:ct-de-compat} and~\ref{lem:de-cuts});
\item the hard direction of CGGR Theorem 1.1 (Theorem~\ref{thm:cggr11-hard}), whose proof
cites Loewy (1986), Theorem 1 (Lemma~\ref{ext:no-cut}) and Hartfiel--Loewy (1984).
\end{itemize}
Lemmas 4.3, 4.4 and 4.5 of the paper and the determinantal representation theorem
(Lemma~\ref{ext:recognition}) are not needed for it.

\section{Conventions and boundaries}

Matrix identities are stated over commutative rings whenever possible, and use arbitrary
finite index types rather than a fixed enumeration. Polynomial support restrictions are
predicates in one ambient multivariate polynomial ring. The rank definition of a cut
uses a field; ranks of polynomial matrices in Lemma 4.3 are taken over a fraction field.
Irreducibility means strong connectivity for nonzero entries, including length-zero paths.
Thus the one-by-one zero matrix is irreducible, and so is the empty matrix. The latter has
$f_A=1$, so the polynomial irreducibility criterion needs a nonempty index type.

A matrix is \emph{symmetrizable} if an invertible diagonal matrix $\mu$ makes $\mu A$
symmetric. Conjugacy to a symmetric matrix can require square roots; the two notions are
only claimed to agree over an algebraically closed field. Diagonal similarity and
diagonal equivalence must be kept apart: the error in the paper's proof is that its second
cut swap is the transpose of the first, hence already diagonally equivalent to it.
Sections 5--7 of the paper beyond the irreducibility criterion and the symmetric-fiber
lemma are outside the current target.
